%%
%% MICRO 2026 Paper: Model-Defined Remapper
%% Based on the official MICRO 2026 LaTeX template (acmart sigconf)
%%
%TC:macro \cite [option:text,text]
%TC:macro \citep [option:text,text]
%TC:macro \citet [option:text,text]
%TC:envir table 0 1
%TC:envir table* 0 1
%TC:envir tabular [ignore] word
%TC:envir displaymath 0 word
%TC:envir math 0 word
%TC:envir comment 0 0

\documentclass[sigconf, screen, review]{acmart}

\AtBeginDocument{%
  \providecommand\BibTeX{{%
    Bib\TeX}}}

\setcopyright{none}
\copyrightyear{2026}
\acmYear{2026}
\acmDOI{XXXXXXX.XXXXXXX}
\acmConference[MICRO 2026]{The 59th IEEE/ACM International Symposium on Microarchitecture}{October 31--November 04, 2026}{Athens, Greece}
\acmISBN{978-X-XXXX-XXXX-X/XX/XX}

%%\acmSubmissionID{123-A56-BU3}

\settopmatter{printfolios=true}
\settopmatter{printacmref=false}

\begin{document}

%% ============================================================
%% TITLE
%% ============================================================
\title{Defect-Tolerant Wafer-Scale NoC Remapping for Disaggregated AI Inference}
\subtitle{\normalsize{MICRO 2026 Submission
    \textbf{\#NaN} -- Confidential Draft -- Do NOT Distribute!!}}

%% ============================================================
%% ABSTRACT (~200 words)
%% ============================================================
\begin{abstract}

% --- Sentence 1-2: Context & problem ---
% Large-scale 2D mesh architectures (wafer-scale engines, chiplet interposers) are
% increasingly adopted for AI and HPC workloads, yet manufacturing defects in cores,
% routers, and links are inevitable at scale.
%
% --- Sentence 3-4: Limitation of existing approaches ---
% Current defect mitigation relies on either hardware redundancy (spare cores/routers)
% that wastes die area, or software discard of irregular regions that underutilizes
% healthy resources. Both approaches expose only a regular sub-topology to software,
% leaving routing resources (e.g., diagonal links) unused.
%
% --- Sentence 5-6: Our approach ---
% We propose Model-Defined Remapper, a framework that decouples the virtual topology
% seen by applications from the defective physical topology. Our approach uses a
% hierarchical spatial search heuristic to map a virtual regular mesh onto all healthy
% nodes of a larger faulty physical mesh, paired with a submesh-hierarchy routing
% scheme that guarantees deadlock-free packet delivery on the resulting irregular topology.
%
% --- Sentence 7-8: Key results ---
% We evaluate on meshes up to 512x512 with realistic fault rates (7\% PE, 3.5\% router,
% 0.7\% link), demonstrating that Model-Defined Remapper achieves [X\%] higher hardware
% utilization than discard-based approaches while maintaining [Y\%] of healthy-mesh
% performance for GEMM and GEMV workloads. We further show that mesh-to-diagonal-mesh
% remapping exploits unused routing resources, reducing average hop distance by [Z\%].

\end{abstract}

\keywords{topology mapping, defect tolerance, network-on-chip, deadlock-free routing, wafer-scale, chiplet}

\maketitle

%% ============================================================
%% 1. INTRODUCTION (1.5--2 pages)
%% ============================================================
\section{Introduction}
\label{sec:intro}

% --- P1: Big picture --- (4-5 sentences)
% The era of large-scale 2D mesh architectures.
% Wafer-scale engines (Cerebras WSE: 850K+ cores), chiplet interposers (AMD MI300),
% and many-core accelerators use 2D mesh/torus as their interconnect topology.
% These architectures are critical for AI training, inference, and HPC workloads
% that demand massive parallelism with structured communication (systolic arrays,
% broadcast-reduce patterns).
% Cite: \cite{lie2019wse, lie2024cerebras, xue2024npu}

% --- P2: The defect problem --- (4-5 sentences)
% At wafer scale, defects are inevitable: PU failures, router faults, link opens/shorts.
% As feature sizes shrink and die areas grow, defect density becomes a first-order concern.
% Cerebras reports ~1% spare cores to compensate \cite{lie2024cerebras}.
% The fundamental tension: applications require regular topologies for structured
% computation (GEMM, GEMV), but physical hardware is irregular due to defects.

% --- P3: Existing approaches and their limitations --- (5-6 sentences)
% Two dominant strategies exist:
% (1) Hardware redundancy: spare cores/routers swap in for defective ones, presenting
%     a regular topology to software (Cerebras WSE, AMD MI300).
%     Limitation: wastes die area on spares; extra routing resources (diagonal links)
%     go unused when only regular mesh is exposed.
% (2) Software discard: find largest defect-free regular sub-mesh, use only that.
%     Limitation: abandons usable-but-irregular regions, low utilization at higher
%     defect rates.
% Neither approach is workload-aware or scales gracefully with defect rate.
% Cite: \cite{zhang2009topology, zhang2010virtualization, cerebras2025sigops}

% --- P4: Our key insight --- (3-4 sentences)
% Insight: we can decouple the virtual topology (what the application programs against)
% from the physical topology (the actual defective hardware).
% Rather than forcing physical regularity, we map a virtual regular mesh onto ALL
% available healthy nodes in a larger faulty physical mesh, with a cost model that
% is aware of the application's communication patterns (GEMM, GEMV).
% The mapping is transparent: applications are unchanged.

% --- P5: Challenges --- (3-4 sentences)
% This approach raises three challenges:
% (C1) The virtual-to-physical embedding problem is NP-hard (generalizes subgraph
%      isomorphism with cost minimization), requiring scalable heuristics.
% (C2) Routing on the resulting irregular physical topology must be deadlock-free,
%      which is non-trivial when faults break the regular structure that standard
%      dimension-order routing relies on.
% (C3) The remapper must scale to large meshes (512x512+) with practical runtime.
% Cite: \cite{rost2020hardness, zhang2009topology}

% --- P6: Our approach --- (3-4 sentences)
% We propose Model-Defined Remapper, which addresses all three challenges:
% A hierarchical spatial search (recursive coordinate bisection + KdTree-indexed
% super nodes) enables scalable mapping with workflow-aware cost optimization.
% A submesh-hierarchy routing scheme (ascend-then-descend, no re-ascent) guarantees
% deadlock-free delivery on arbitrarily irregular topologies.
% We further demonstrate that remapping a virtual 4-connected mesh onto a physical
% 8-connected (diagonal) mesh exploits otherwise-unused routing resources.

% --- P7: Contributions --- (bulleted list)
% We make the following contributions:
\begin{itemize}
  \item We formalize the defect-tolerant virtual-to-physical topology remapping
        problem for large-scale 2D mesh architectures and show it is NP-hard.
  \item We propose a hierarchical, workflow-aware mapping heuristic that scales to
        512$\times$512 meshes, using KdTree-indexed spatial search with parallel
        candidate evaluation.
  \item We design a submesh-hierarchy routing scheme that guarantees deadlock-free
        packet delivery on arbitrarily defective topologies, for both 4-connected
        and 8-connected meshes.
  \item We demonstrate that mesh-to-diagonal-mesh remapping exploits unused routing
        resources (e.g., diagonal links in Cerebras-style fabrics), reducing hop
        distance and congestion.
  \item We evaluate across mesh sizes up to 512$\times$512, multiple fault rates, and
        GEMM/GEMV workloads, showing [X\%] utilization improvement over discard-based
        approaches with [Y\%] of healthy-mesh performance.
\end{itemize}

%% ============================================================
%% 2. BACKGROUND AND MOTIVATION (1--1.5 pages)
%% ============================================================
\section{Background and Motivation}
\label{sec:background}

% ---------- 2.1 ----------
\subsection{Large-Scale 2D Mesh Architectures}
\label{sec:bg-arch}

% --- P1: Why 2D mesh dominates --- (4-5 sentences)
% 2D mesh is the dominant on-chip topology for large-scale architectures because:
% (1) Manufacturing-friendly: planar layout maps directly to silicon/interposer substrate.
% (2) Scalable: O(N) links, constant router radix.
% (3) Enables structured computation: systolic arrays (GEMM), broadcast-reduce (GEMV).
% Examples: Cerebras WSE (2D mesh, 850K cores), Google TPU (2D torus),
%           Intel Ponte Vecchio (mesh of Xe cores).

% --- P2: Why not other topologies --- (3-4 sentences)
% Clos/fat-tree: easier defect tolerance (stage-level swap), but multi-stage switching
% doesn't map to 2D silicon layout; suited for switches/routers, not on-chip.
% Fully-connected: trivial routing but O(N^2) links, unaffordable at scale.
% 2D mesh/torus is the practical choice for 100+ core chips.
% Cite: \cite{lie2019wse, yin2023modular}

% ---------- 2.2 ----------
\subsection{Defects and Current Mitigation}
\label{sec:bg-defect}

% --- P1: Defect types --- (3-4 sentences)
% Three types of permanent manufacturing defects in 2D mesh architectures:
% (1) Processing Element (PE/core) faults: stuck-at, logic errors.
% (2) Router faults: crossbar or buffer failures.
% (3) Link faults: open/short circuits in inter-router wires.
% At wafer scale, defect rates of 1-10% per component type are realistic.

% --- P2: Hardware redundancy approach --- (4-5 sentences)
% Approach 1: Design with spare cores, routers, and links.
% Cerebras WSE: ~1% spare cores, redundant fabric links; distributed autonomous
% repair logic restores a regular 2D mesh topology.
% AMD MI300: redundant XCDs (chiplet-level redundancy).
% Key limitation: spares consume die area/power; diagonal links in physical fabric
% are unused when only regular 4-connected mesh is exposed to software.
% Cite: \cite{lie2024cerebras, cerebras2025sigops}

% --- P3: Software discard approach --- (3-4 sentences)
% Approach 2: Identify largest defect-free rectangular sub-mesh; use only that.
% Simple but wasteful: all healthy nodes outside the sub-mesh are abandoned.
% Utilization degrades rapidly with defect rate.
% Prior work: Zhang et al. proved topology reconfiguration is NP-complete \cite{zhang2009topology}.

% ---------- 2.3 ----------
\subsection{The Utilization Gap}
\label{sec:bg-gap}

% --- Motivating data/figure --- (4-5 sentences + figure)
% [FIGURE 1: Motivating example]
% Left: 12x12 physical mesh with 7% PE faults, 3.5% router faults.
% Middle: Largest regular sub-mesh (discard approach) — show wasted healthy nodes.
% Right: Our remapping — virtual 8x8 mesh mapped onto all healthy nodes.
% Quantify: discard wastes X% of healthy resources; our approach uses Y%.
%
% Key observation: even with Cerebras-style diagonal links in the physical fabric,
% the discard/redundancy approach only exposes a 4-connected regular mesh,
% leaving diagonal routing resources unused.

% ---------- 2.4 ----------
\subsection{Regularity Requirements for Structured Computation}
\label{sec:bg-regularity}

% --- P1: Why applications need regular topology --- (4-5 sentences)
% Structured kernels (GEMM, GEMV, convolution, FFT) rely on regular communication patterns:
% Systolic GEMM: data flows along rows and columns in lockstep.
% GEMV: broadcast input vector along rows, reduce partial sums along columns.
% These patterns assume uniform neighbor connectivity and predictable hop latency
% that dimension-order routing on a regular mesh provides.
% Irregular topology breaks these assumptions: non-uniform hop counts cause
% synchronization overhead and load imbalance.

% --- P2: Our insight --- (2-3 sentences)
% Key insight: the application only needs to *see* a regular topology.
% If we can transparently remap the virtual regular mesh to the physical irregular mesh
% while preserving communication semantics, applications need not change.
% The cost is non-uniform physical hop distances, which our cost model minimizes.

%% ============================================================
%% 3. MODEL-DEFINED REMAPPER (3--4 pages)
%% ============================================================
\section{Model-Defined Remapper}
\label{sec:design}

% ---------- 3.1 ----------
\subsection{Overview}
\label{sec:design-overview}

% [FIGURE 2: System overview / pipeline diagram]
% Input:  Virtual mesh spec (W x H) + Physical mesh with fault map + Workload (GEMM/GEMV)
% Stage 1: Fault-aware partitioning (RCB -> submeshes -> super nodes)
% Stage 2: Hierarchical mapping (anchor init -> frontier expansion -> cost evaluation)
% Stage 3: Routing construction (submesh-hierarchy paths -> routing database)
% Output: Coordinate map (v_coord -> p_coord) + Deadlock-free routing tables

% --- P1: Design philosophy --- (3-4 sentences)
% Model-Defined Remapper takes a "virtual topology first" approach:
% the application defines its desired regular mesh; the remapper finds the best
% physical embedding. The mapping is one-time (post-manufacturing / boot-time),
% not runtime. The framework is topology-parametric: supports both 4-connected
% and 8-connected physical meshes via pluggable distance metrics and routers.

% ---------- 3.2 ----------
\subsection{Problem Formulation}
\label{sec:formulation}

% --- P1: Formal definition --- (5-6 sentences + math)
% Given:
%   - Virtual mesh G_v = (V_v, E_v): W x H regular 2D mesh
%   - Physical mesh G_p = (V_p, E_p): W' x H' mesh with faulty nodes/edges removed
%     (W' > W, H' > H to accommodate faults)
%   - Workload W: set of data flows (unicast, multicast, reduce, alltoall)
% Find: injection f: V_v -> V_p such that:
%   (1) f is injective (one-to-one)
%   (2) For every virtual edge (u,v) in E_v, there exists a physical path from f(u) to f(v) in G_p
%   (3) Cost(f, W) is minimized, where Cost accounts for:
%       - Hop stretch: physical distance vs. virtual distance for each mapped edge
%       - Congestion: link utilization based on workload data flows
%       - Workflow cost: weighted by communication intensity

% --- P2: Complexity --- (3-4 sentences)
% Theorem: The mapping problem is NP-hard.
% Proof sketch: Reduction from subgraph isomorphism (known NP-complete).
% Our problem is strictly harder: we minimize a cost function over valid embeddings
% (optimization variant), which is inapproximable in general \cite{rost2020hardness}.
% This motivates our heuristic approach.

% ---------- 3.3 ----------
\subsection{Fault-Aware Mesh Partitioning}
\label{sec:partitioning}

% [FIGURE 3: Partitioning example]
% Show a faulty physical mesh being recursively bisected into submeshes,
% then submeshes grouped into super nodes.

% --- P1: Recursive Coordinate Bisection (RCB) --- (4-5 sentences)
% We partition the physical mesh into rectangular submeshes via RCB:
% Recursively bisect along the longer dimension, splitting at midpoint.
% Faulty nodes are excluded from submesh interiors.
% Each submesh has: bounding box (ll, ur), center, set of healthy nodes.
% Controlled by max_block_size parameter (e.g., 4).

% --- P2: Submesh adjacency graph --- (3-4 sentences)
% Build an adjacency graph over submeshes: two submeshes are adjacent if they
% share a physical boundary with at least one healthy link crossing it.
% This graph is the foundation for both candidate search and routing.

% --- P3: Super node construction --- (3-4 sentences)
% Group nearby submeshes into super nodes (BFS within super_node_range hops
% on the submesh adjacency graph).
% Build a KdTree over super node centers for O(log m) nearest-neighbor queries.
% This two-level hierarchy (submeshes within super nodes) enables efficient
% candidate search on meshes up to 512x512+.

% ---------- 3.4 ----------
\subsection{Hierarchical Mapping Algorithm}
\label{sec:algorithm}

% [ALGORITHM 1: Pseudocode for mapping pipeline]
% function MAP(G_v, G_p, W):
%   submeshes, adj, super_nodes <- PARTITION(G_p)
%   anchor <- SELECT_ANCHOR(G_v, submeshes)
%   mapping <- INIT_MAPPING(anchor)
%   while mapping is incomplete:
%     mapping <- EXPAND_FRONTIER(mapping, super_nodes, W)
%   routing <- BUILD_ROUTING(mapping, submeshes, adj)
%   return mapping, routing

% --- P1: Anchor initialization --- (3-4 sentences)
% Seed the mapping by placing the virtual mesh center at the physical mesh's
% root submesh center (closest to geometric center).
% Map the anchor's immediate neighbors to establish initial cost baseline.
% The anchor determines the global orientation of the mapping.

% --- P2: Frontier expansion --- (5-6 sentences)
% [ALGORITHM 2: Pseudocode for EXPAND_FRONTIER]
% For each unmapped virtual node adjacent to the current frontier:
%   (1) Compute ideal physical location via geometric transform from existing mapping.
%   (2) Level-1 search: KdTree query for nearest search_nearest_n super nodes.
%   (3) Level-2 search: Score submesh centers within top super nodes; select top-K.
%   (4) Parallel evaluation: For each of expansion_trials physical candidates,
%       compute cost = node_cost + hop_cost_weight * path_cost + workflow_cost_weight * workflow_cost.
%   (5) Commit best candidate; update frontier and candidate pool.
% This two-level search avoids O(m) exhaustive scan over the entire physical mesh.

% --- P3: Cost model --- (5-6 sentences)
% The cost function has three components:
% (1) Coordinate cost: penalizes stretch between mapped virtual neighbors.
%     Uses 2-hop lookahead with exponential decay weighting.
%     Manhattan distance for 4-connected; Chebyshev for 8-connected.
% (2) Congestion cost: tracks per-link utilization via LinkUsage.
%     Penalizes placing nodes that would overload already-congested paths.
% (3) Workflow cost: weights by communication intensity from workload analysis.
%     GEMM systolic flows are weighted differently than GEMV broadcast-reduce.
% RANSAC geometric consistency check filters outlier placements.

% ---------- 3.5 ----------
\subsection{Deadlock-Free Routing}
\label{sec:routing}

% [FIGURE 4: Submesh hierarchy routing example]
% Show a packet path: ascend through submesh hierarchy to common ancestor,
% then descend to destination submesh. Annotate "no re-ascent" constraint.

% --- P1: The deadlock challenge --- (3-4 sentences)
% On a regular mesh, XY dimension-order routing is trivially deadlock-free.
% After remapping, physical paths between virtual neighbors may traverse multiple
% submeshes in complex patterns. Naive shortest-path routing on the irregular
% topology can create cyclic channel dependencies, leading to deadlock.
% Cite: \cite{dally1987deadlock, duato1993theory}

% --- P2: Submesh hierarchy routing --- (5-6 sentences)
% We construct a routing hierarchy from the submesh adjacency tree (rooted BFS tree).
% For any src-dst pair, the physical path is computed in two phases:
%   Phase 1 (ascend): route from source submesh up to the lowest common ancestor (LCA)
%   in the submesh tree.
%   Phase 2 (descend): route from LCA down to destination submesh.
% Constraint: once a packet begins descending, it never re-ascends.
% This ascend-then-descend (up*/down*-style) constraint creates an acyclic channel
% dependency graph, guaranteeing deadlock freedom by the Dally-Seitz theorem.
% Within each submesh, standard XY routing (or diagonal XY) is used.
% Cite: \cite{dally1987deadlock, ariadne2011}

% --- P3: Routing database --- (3-4 sentences)
% All physical paths are precomputed and cached in a RoutingDatabase.
% For multicast/broadcast data flows, paths are merged into BFS routing trees
% to minimize duplicate transmissions.
% The routing map is constructed once at mapping time and stored in routing tables.

% --- P4: Correctness argument --- (3-4 sentences)
% Theorem: The submesh-hierarchy routing scheme is deadlock-free.
% Proof sketch: The ascend-then-descend constraint induces a total order on channel
% usage (channels used in ascending phase < channels in descending phase).
% By the Dally-Seitz CDG acyclicity condition, no deadlock can occur.
% This holds for arbitrary fault patterns as long as the submesh adjacency graph
% is connected (which our partitioner guarantees for viable mappings).

% ---------- 3.6 ----------
\subsection{Exploiting Heterogeneous Physical Links}
\label{sec:diagonal}

% --- P1: Mesh-to-diagonal-mesh remapping --- (4-5 sentences)
% Modern large-scale chips (e.g., Cerebras WSE) have physical fabrics with diagonal
% links for redundancy, but expose only a regular 4-connected mesh to software.
% Our remapper can map a virtual 4-connected mesh onto an 8-connected (diagonal)
% physical mesh, leveraging diagonal links for shorter physical paths.
% The cost model switches to Chebyshev distance (max(dx,dy)) for 8-connected,
% naturally preferring diagonal shortcuts when available.
% This reduces average hop count and alleviates congestion on critical paths.

%% ============================================================
%% 4. IMPLEMENTATION (0.5--1 page)
%% ============================================================
\section{Implementation}
\label{sec:implementation}

% --- P1: Remapper implementation --- (4-5 sentences)
% Model-Defined Remapper is implemented in Rust (~15K LoC) as a modular workspace:
% NoC module (topology models), IR module (workload representation),
% Mapper module (partitioning, mapping, routing), Simulator module (performance prediction).
% Parallel candidate evaluation uses Rayon (work-stealing thread pool).
% Seeded ChaCha20 RNG ensures reproducibility across runs.

% --- P2: Simulator --- (3-4 sentences)
% We use a custom cycle-level NoC simulator supporting wormhole routing with
% virtual channels, configurable buffer sizes, and bandwidth parameters.
% A lightweight flat-array simulator variant handles 512x512+ meshes efficiently.
% Performance prediction uses a communication predictor (wormhole pipelining model)
% and a computation predictor (per-PE step model) with three congestion modes
% (best/average/worst).

% --- P3: Workload modeling --- (3-4 sentences)
% Workloads are modeled as WorkGraphs: directed acyclic graphs of communication
% and computation operations.
% We implement two canonical workloads:
% (1) Systolic GEMM: row-column data flows with step-synchronized compute.
% (2) Broadcast-reduce GEMV: broadcast input vector, local multiply, reduce partial sums.
% Both use the IR's DataFlow abstraction (unicast, multicast, reduce, alltoall).

% [TABLE 1: Implementation parameters]
% bandwidth_per_cycle: 32 bits/cycle
% flops_per_pu_per_cycle: 16
% hop_latency_cycles: 1.0
% buffer_size: 4 (VC buffer depth)
% mapper parameters: max_block_size=4, search_nearest_n=5, search_top_k=3, expansion_trials=10

%% ============================================================
%% 5. EVALUATION (2.5--3 pages)
%% ============================================================
\section{Evaluation}
\label{sec:evaluation}

% We answer the following questions:
% Q1: How much utilization does Model-Defined Remapper recover vs. discard-based approaches?
% Q2: What is the performance overhead of remapped execution vs. healthy mesh?
% Q3: Does mesh-to-diagonal-mesh remapping improve performance?
% Q4: How does the approach scale with mesh size and defect rate?
% Q5: What is the mapping runtime?

% ---------- 5.1 ----------
\subsection{Methodology}
\label{sec:methodology}

% --- P1: Experimental setup --- (4-5 sentences)
% Virtual mesh sizes: 64x64, 128x128, 256x256, 512x512.
% Physical mesh: 1.5x virtual size (e.g., 96x96, 192x192, 384x384, 768x768).
% Fault injection: seeded random (ChaCha20, seed=42).
%   PE fault rate: 7%, Router fault rate: 3.5%, Link fault rate: 0.7%.
% Workloads: GEMM (matrix sizes 128, 256, 512), GEMV (same).

% --- P2: Baselines --- (3-4 sentences)
% (1) Healthy: ideal performance on defect-free mesh (upper bound).
% (2) Discard: largest defect-free rectangular sub-mesh (lower bound on utilization).
% (3) Zhang et al. (TVLSI 2009): simulated annealing topology reconfiguration (prior SOTA).
% (4) ARIADNE (PACT 2011): topology-agnostic up*/down* reconfiguration (deadlock-free baseline).

% --- P3: Metrics --- (3-4 sentences)
% Utilization: fraction of healthy nodes used vs. total healthy nodes available.
% Performance: predicted execution cycles (best/average/worst congestion).
% Mapping quality: node cost (stretch), path cost (congestion).
% Mapping time: wall-clock seconds for the mapping algorithm.

% [TABLE 2: Experimental configuration summary]

% ---------- 5.2 ----------
\subsection{Utilization Recovery}
\label{sec:results-util}

% [FIGURE 5: Utilization comparison bar chart]
% Compare utilization across: Discard, Zhang et al., ARIADNE, Ours
% Across mesh sizes 64x64 to 512x512.
%
% --- P1: Main result --- (4-5 sentences)
% Model-Defined Remapper achieves [X%] utilization vs. [Y%] for discard
% and [Z%] for Zhang et al. at 512x512 scale.
% The gap widens at higher defect rates because our hierarchical search
% can find pockets of healthy nodes that discard misses.
% At 7% PE fault rate on 512x512, discard wastes ~[A]% of healthy nodes;
% our approach wastes only ~[B]%.

% ---------- 5.3 ----------
\subsection{Performance}
\label{sec:results-perf}

% [FIGURE 6: Performance (cycles) comparison — GEMM]
% [FIGURE 7: Performance (cycles) comparison — GEMV]
% Lines: Healthy, Ours (best/avg/worst), Zhang, ARIADNE, Discard
% X-axis: mesh size. Y-axis: execution cycles.
%
% --- P1: GEMM results --- (4-5 sentences)
% For systolic GEMM, our remapper achieves [X%] of healthy-mesh performance
% at 256x256, compared to [Y%] for Zhang and [Z%] for ARIADNE.
% The workflow-aware cost model places communicating virtual neighbors closer
% together on critical systolic paths.

% --- P2: GEMV results --- (3-4 sentences)
% For broadcast-reduce GEMV, the performance gap is [smaller/larger] because
% [broadcast patterns are more/less sensitive to hop stretch].
% Our approach outperforms baselines by [X%] due to congestion-aware link placement.

% ---------- 5.4 ----------
\subsection{Diagonal Link Exploitation}
\label{sec:results-diagonal}

% [FIGURE 8: Mesh-to-mesh vs. mesh-to-diagonal-mesh comparison]
% Compare: physical 4-connected (standard remapping) vs. physical 8-connected
% (diagonal mesh remapping).
% Metrics: average hop distance, congestion, execution cycles.
%
% --- P1: Hop reduction --- (3-4 sentences)
% Remapping onto 8-connected physical mesh reduces average hop distance by [X%]
% compared to 4-connected physical mesh.
% Diagonal links provide shortcuts for non-axis-aligned virtual neighbors.

% --- P2: Performance improvement --- (3-4 sentences)
% This translates to [Y%] cycle reduction for GEMM and [Z%] for GEMV.
% The improvement demonstrates that exposing all physical routing resources
% to the remapper — rather than hiding them behind a regular 4-connected view —
% yields tangible performance gains.

% ---------- 5.5 ----------
\subsection{Scalability}
\label{sec:scalability}

% [FIGURE 9: Mapping time vs. mesh size (log-log plot)]
% [TABLE 3: Mapping statistics at each scale]
%
% --- P1: Runtime scaling --- (4-5 sentences)
% Mapping time scales as O(N * log M) where N = virtual mesh size, M = physical mesh size.
% At 512x512 (virtual) -> 768x768 (physical), mapping completes in [X] seconds
% on 8 threads.
% The hierarchical KdTree search keeps per-node cost at O(log M),
% avoiding the O(M) exhaustive scan that limits prior approaches.
% Zhang et al.'s SA-based approach does not scale beyond ~128x128.

% --- P2: Mapping quality vs. scale --- (3-4 sentences)
% Node cost and path cost remain bounded as mesh size increases,
% demonstrating that our hierarchical partitioning maintains mapping quality.
% Number of submeshes grows linearly with physical mesh area.

% ---------- 5.6 ----------
\subsection{Sensitivity Analysis}
\label{sec:sensitivity}

% [FIGURE 10: Performance vs. defect rate (sweep 1%--15%)]
% [FIGURE 11: Performance vs. defective_mesh_factor (1.2x, 1.5x, 2.0x)]
%
% --- P1: Defect rate sensitivity --- (3-4 sentences)
% Performance degrades gracefully as defect rate increases from 1% to 15%.
% At 15% PE fault rate, our approach still achieves [X%] of healthy performance,
% while discard utilization drops below [Y%].

% --- P2: Physical mesh sizing --- (3-4 sentences)
% Larger defective_mesh_factor (more physical area per virtual node) improves
% mapping quality but increases physical area.
% Factor 1.5x provides the best quality/area tradeoff for 7% defect rates.

% --- P3: Mapper parameter sensitivity --- (3-4 sentences)
% search_nearest_n and expansion_trials: diminishing returns beyond 5 and 10 respectively.
% hop_cost_weight: workflow-aware weighting improves GEMM by [X%] over topology-only cost.

%% ============================================================
%% 6. RELATED WORK (0.5--1 page)
%% ============================================================
\section{Related Work}
\label{sec:related}

% ---------- 6.1 ----------
\subsection{Topology Reconfiguration for Defect Tolerance}
% Zhang et al. (DATE 2008, TVLSI 2009) proposed core-level redundancy with topology
% reconfiguration via simulated annealing, proved NP-completeness.
% Zhang et al. (DATE 2010) extended with performance-asymmetry awareness.
% Ke et al. (ASP-DAC 2012) used Hungarian algorithm for timing-similarity mapping.
% Qian et al. (JPDC 2024) proposed greedy-memetic algorithm.
% Ding et al. (TACO 2025) introduced two-stage degradation-based reconfiguration.
% Key difference: all prior works evaluate at <128 cores, use topology-only cost
% (no workload awareness), and don't address deadlock-free routing explicitly.
% Cite: \cite{zhang2009topology, zhang2010virtualization, ke2012hungarian, qian2024memetic, ding2025twostage}

% ---------- 6.2 ----------
\subsection{Deadlock-Free Routing on Irregular Topologies}
% Foundational: Dally-Seitz CDG theorem, Turn Model (Glass \& Ni), Duato's theory.
% ARIADNE (PACT 2011): topology-agnostic reconfiguration with up*/down* routing.
% uDIREC (MICRO 2013): unified diagnosis + MOUNT routing on faulty NoC.
% Vicis (DATE 2009): resilient routing tolerating 50% router failures.
% DeFT (DATE 2022): deadlock-free fault-tolerant routing for 2.5D chiplets.
% Key difference: these works handle routing but don't provide virtual topology
% abstraction or workload-aware mapping.
% Cite: \cite{dally1987deadlock, ariadne2011, udirec2013, vicis2009, deft2022}

% ---------- 6.3 ----------
\subsection{NoC Application Mapping}
% Hu \& Marculescu (TCAD 2005): branch-and-bound energy-aware mapping.
% Kadri \& Koudil (JSA 2019): survey of fault-tolerant mapping techniques.
% Neurocomputing 2020: deadlock-free mapping for neural network chips.
% Key difference: task mapping assigns tasks to PEs but doesn't remap the
% entire virtual topology; most works assume healthy topology.
% Cite: \cite{hu2005mapping, kadri2019survey}

% ---------- 6.4 ----------
\subsection{Accelerator Virtualization}
% vNPU (ISCA 2025): topology-aware virtualization for inter-core connected NPUs.
% Xue et al. (MICRO 2024): hardware-assisted NPU virtualization for cloud.
% Cerebras (SIGOPS 2025): WSE driver remaps around defects.
% Key difference: vNPU targets NPU multi-tenancy (not defect tolerance);
% Cerebras only restores regularity (doesn't optimize for workload or exploit extra links).
% Cite: \cite{feng2025vnpu, xue2024npu, cerebras2025sigops}

%% ============================================================
%% 7. CONCLUSION (0.5 page)
%% ============================================================
\section{Conclusion}
\label{sec:conclusion}

% --- P1: Summary --- (4-5 sentences)
% We presented Model-Defined Remapper, a framework for defect-tolerant
% virtual-to-physical topology mapping on large-scale 2D mesh architectures.
% Our hierarchical mapping heuristic scales to 512x512 meshes while optimizing
% for workload-specific communication patterns.
% Our submesh-hierarchy routing scheme guarantees deadlock-free delivery on
% arbitrarily defective topologies.
% We demonstrated that mesh-to-diagonal-mesh remapping exploits unused routing
% resources in Cerebras-style fabrics.

% --- P2: Key results --- (2-3 sentences)
% Evaluation across multiple scales, fault rates, and workloads shows [X%]
% utilization recovery over discard-based approaches and [Y%] of healthy-mesh
% performance for GEMM and GEMV workloads.
% Our approach is the first to combine large-scale topology remapping,
% workload-aware optimization, and provably deadlock-free routing in a unified framework.

% --- P3: Future directions --- (2-3 sentences)
% Future work includes runtime remapping for online fault discovery,
% extension to non-mesh virtual topologies (torus, 3D mesh),
% and integration with real hardware platforms (Cerebras SDK, chiplet interposers).

%% ============================================================
%% REFERENCES
%% ============================================================
\bibliographystyle{ACM-Reference-Format}
\bibliography{references}

\end{document}
